%============================================================================
% PREÁMBULO
%============================================================================
\documentclass[letterpaper,12pt]{report}
\usepackage[a4paper, margin=2cm, marginparwidth=3cm]{geometry} % Configura el tamaño del papel y los márgenes del documento
\usepackage[spanish]{babel}
\selectlanguage{spanish}
\usepackage[utf8]{inputenc}
\usepackage[numbers]{natbib}
\usepackage{graphicx,setspace,longtable,pdflscape,color}
\usepackage[colorlinks=true,linkcolor=blue,
           urlcolor=blue,citecolor=blue]{hyperref}
\usepackage{titlesec,url}
\usepackage[nottoc,notlot,notlof]{tocbibind}
%\usepackage[margin=2.54 cm]{geometry}
\usepackage[version=4]{mhchem}
\usepackage{amssymb, amsmath} 
\usepackage{float}
\usepackage{caption}
\usepackage{subcaption}
\usepackage{xspace}
\usepackage{fancyhdr}
\usepackage{fontawesome5}
\usepackage{amsmath}
\usepackage{physics}
\usepackage{tikz}
\usetikzlibrary{patterns, decorations.pathreplacing}
\usepackage[most]{tcolorbox}

\newtcolorbox{infobox}[2][blue]{
    colback=#1!5!white,
    colframe=#1!75!black,
    title={#2},
    fonttitle=\bfseries,
    coltitle=black,
    boxrule=0.8pt,
    arc=4pt,
    left=6pt,
    right=6pt,
    top=6pt,
    bottom=6pt,
    breakable
}

\usepackage{cancel}
\newcommand{\colorcancel}[2]{\textcolor{#1}{\cancel{\textcolor{black}{#2}}}}

% Comando personalizado para CosmoLattice
\newcommand{\CosmoLattice}{\ensuremath{\mathcal{C}}\text{osmo}\ensuremath{\mathcal{L}}\text{attice}\xspace}

% Definición del estilo para el cuerpo de la tesis (pág xx de yy)
\fancypagestyle{estilotesis}{
  \fancyhf{} 
  \cfoot{pág \thepage\ de \pageref*{ultimaPagina}} 
  \renewcommand{\headrulewidth}{0pt}
}

\definecolor{gold}{RGB}{212, 175, 55}     % Oro metálico
\definecolor{silver}{RGB}{125, 125, 125}  % Plata clásico
\definecolor{ForestGreen}{HTML}{14BA24}
\definecolor{esmerald}{RGB}{80, 200, 120}
\definecolor{Green}{rgb}{0.13, 0.55, 0.13}

% ============================================================================
% CONFIGURACIÓN DE ESPACIADO EN ECUACIONES
% ============================================================================
\setlength{\abovedisplayskip}{6pt plus 2pt minus 2pt}
\setlength{\belowdisplayskip}{6pt plus 2pt minus 2pt}
\setlength{\abovedisplayshortskip}{0pt plus 2pt}
\setlength{\belowdisplayshortskip}{6pt plus 2pt minus 2pt}

\usepackage{tocloft}

% Ajusta el espacio ANTES de "Índice general" (podés usar valores negativos)
\setlength{\cftbeforetoctitleskip}{0cm} 

% Ajusta el espacio DESPUÉS de "Índice general" y antes del primer ítem
\setlength{\cftaftertoctitleskip}{0cm}

%============================================================================
% COMIENZO DEL DOCUMENTO
%============================================================================
\begin{document}
  \section{Presentación}
    Presentación

  \section{Hoja de ruta}
    Esta charla la vamos a dividir en 4 partes: por un lado empezaré contando el marco teórico en el que está metido esta charla, luego tendremos dos secciones para contar lo que estuvimos trabajando y finalmente qué se concluye.

  \section{Título Intro}
    Empecemos con la intro

  \section*{Indice Intro}
    La misma vamos a separarla en dos partes: por un lado contar el contexto cosmológico y por el otro lado qué son las ondas gravitacionales
  
  \section{Historia del Universo}
    Para empezar con la parte de Cosmología, empecemos marcando que en el modelo cosmológico estándar, el universo nace con el Big Bang, seguido por una etapa conocida como Inflación, un Recalentamiento, se crean los primeros átomos en BBN, luego la luz deja de colisionar con las partículas en el plasma primordial permitiendo viajar hacia nosotros y formar el CMB, se forman las primeras estrellas y galaxias hasta llegar a hoy. En particular, a nosotros nos va a interesar estudiar la región en verde que comprende al final de la Inflación y principio del Recalentamiento
  
  \section{Inflación} 
    La etapa inflacionaria es un período en el cual el universo se expande de forma acelerada permitiendo resolver algunos problemas con las observaciones actuales tales como el problema del horizonte, el problema de la planitud y la sobrepoblación de defectos topológicos.

    Este período típicamente se modela a partir de la incorporación de un único campo escalar conocido como inflatón y cuya dinámica está descripta por la siguiente acción, que a partir de derivar las ecucaciones de movimiento se puede obtener el siguiente tensor de energía-momento, el cual permite definir una densidad de energía y presión asociada.

    A partir de observar esta densidad de energía y presión, se pede observar que si la energía potencial del campo es mucho mayor que la cinética, entonces la presión es negativa, dando lugar a una expansión acelerada del universo.
  
  \section{slow-roll}
    La aproximación de tomar a la energía potencial como mucho mayor a la cinética se conoce como condición de \textit{slow-roll} o rodadura lenta. Esto se debe a que uno piensa a los potenciales del inflatón como lo que se observa acá abajo, en donde se tiene una región en que el inflatón va a rodar suavemente sobre su energía potencial hasta llegar al punto en que las dos energías son comparables, punto el cual se rompe esta aproximación y se considera como el fin de la Inflación y el inicio de la etapa conocida como Recalentamiento.
    
    Para caracterizar esta condición de \textit{slow-roll} y poder marcar un punto en el cual se da esta ruptura se definen a los parámetros de \textit{slow-roll} de esta forma y cuando se rompe la aproximación se obtiene que estos son cercanos a $1$. El calculo de estos parámetros, por tanto, resultarán de interés siendo que darán las condiciones iniciales para lo que charlaremos más adelante.

  \section{Perturbaciones inflación}
    Finalmente, para comprender el estudio de las perturbaciones y en particular de las ondas gravitacionales, vamos a querer entender cómo influye la inflación en las inhomogeneidades del universo. La idea, entonces es que cuando tenemos la expansión violenta, las fluctuaciones cuánticas del vacio se amplifican, saliendo del horizonte de Hubble y congelandose. Esta idea permite servir como semilla para las anisotropías en el CMB y la formación de estructuras.

    A su vez, estas perturbaciones se pueden separar en dos: perturbaciones escalares y tensoriales que van a estar desciptas por los índices espectrales y las amplitudes. En particular nos interesará $r$ dado que dá idea de la amplitud relativa de las GWs con respecto a las perturbaciones de densidad del fluido.

  \section{Recalentamiento}
    La época de Recalentamiento permite hacer de transición entre la Inflación y el momento en que el plasma está totalmente termalizado permitiendo la nucleosíntesis primordial. A esta temperatura a la que ocurre la termalización la conocemos como $T_{reh}$ y existe un amplio rengo de valores que puede tomar la misma, siendo estas cotas las que se pueden ver acá.

    Esta etapa puede ser dividida en una primera parte conocida como \textit{Preheating} que es la que nos interesará estudiar y en la cual se producen partículas y ondas gravitacionales, una época turbulenta y una época de equilibrio.

  \section{Preheating}
    Durante \textit{Preheating}, el inflatón oscila alrededor del mínimo de potencial, este decaimiento puede ser modelado de distintas maneras entre las que se resaltan por el interés de este trabajo los modelos monomiales, los T-models y los E-models. Estas oscilaciones, cuando no se consideran efectos no lineales se consideran coherentes. Mientras esto ocurra, el universo se puede modelar a partir de una dominancia de radiación, materia u otro tipo de materia, siendo la ecuación $\omega_{eff}$ la que permite para modelos monomiales entender esta relación. En esta ecuación se puede observar que si tenemos un potencial cuadrático tendremos una ecuación de estado que caracteriza a la materia, mientras que si tenemos un potencial cuártico tendremos una ecuación de estado que caracteriza a la radiación.

  \section{omega efectivo}
    Lo que mencionamos en la diapo anterior representa a lo que ocurre para los primeros momentos de \textit{Preheating}, pero a medida que aumenta el tiempo, como se observa en la Figura, la ecuación de estado tiende a 1/3 para poder acoplarse con la etapa siguiente. Esto ocurre con excepción del potencial cuadrático para el cual la ecuación de estado efectiva a tiempos largos se mantiene como de tipo materia.

  \section{GWs}
    Las ondas gravitacionales aparecen en la Relatividad General como soluciones a las ecuaciones de Einstein linealizadas. Estas van a describir perturbaciones tensoriales de la métrica escritas como ... Siendo $\bar{g}$ la métrica que describe al fondo y $h$ la parte que describe a las perturbaciones y por tanto a las GWs 

    Típicamente estas perturbaciones se estudian haciendo uso del gauge TT que permite escribir a las perturbaciones a partir de dos polarizaciones conocidas como $h_+$ y $h_\times$ y que para un conjunto de puntos en el espacio estas generan este efecto. Esta elección de gauge considera a la perturbación como transversal a la dirección de propagación y sin traza.

  \section{GW en Cosmo}
    Al utilizar una métrica de fondo de tipo cosmológico se  puede describir a la distibución de energía-momento para las ondas como ... 

    Con lo que podemos definir a la densidad de energía fraccionaria de las ondas gravitacionales observando la primera componente del tensor.
    
    A su vez, las ondas pueden escribirse en términos de sus polarizaciones como ... siendo $h_r$ las polarizaciones y $e_{ij}^r$ el un proyector para cada una de estas polarizaciones.

  \section{Ecuacion de onda Cosmo}
    A partir de la descomposición en las polarizaciones para las ondas, se puede obtener una ecuación de ondas utilizando el fondo de tipo FLRW y nos queda algo de la forma... donde $\Pi$ describe a las perturbaciones tensoriales del tensior de energía-impulso que se puede escribir sacando la parte de fondo del tensor de energía-momento. A su vez, el calculo de este tensor permite definir las proyecciones en polarizaciones y en el gauge TT de la siguiente forma...

  \section{Fondos estocásticos}
    La observación de estas ondas se puede medir a partir de señales coherentes formadas por un único evento o a partir de SGWB que se forman como suma incoherente de ondas tal como se producen ante procesos globales y no aislados.

    Los procesos cosmológicos que forman ondas gravitacionales forman fondos por lo que su detección deberá hacerse a partir de buscar correlaciones entre la señal medida por varios detectores o muchos puntos en el espacio-tiempo.

    Estos fondos se esperan que sean gaussianos, isótropos, homogéneos y no polarizados por lo que su información quedaría totalmente codificada en la función de dos puntos dónde $h_c$ representa la deformación del espacio-tiempo y se relaciona con la densidad de ondas gravitacionales a partir de ...
    
  \section{Detección de SGWB}
    Como se mencionó en la diapo anterior, la detección de estos fondos se realiza a partir de buscar correlaciones en las señales de múltiples detectores, esto podemos pensarlo a partir de pensar que la señal que se mide es parte la onda esperada y parte el ruido instrumental, por lo que si se asume que las señales de las ondas no están correlacionadas con el ruido junto con que el ruido no está correlacionado entre sí, se tiene que la correlación medida es la correlación entre las señales de la onda.

  \section{Fuentes y detectores de SGWB}
    En la figura armamos un compendio de distintas fuentes de ondas gravitacionales y detectores que están actualmente en uso (marcados con línea punteada) y futuros (marcados con línea continua). Los detectores actualmente en uso son Planck/Byceps que utilizan los modos B del CMB para intentar medir las GWs formadas por las fluctuaciones cuánticas, LIGO/VIRGO/KAGRA funcionan como interferómetros tipo Michelson en los cuales al medir una onda gravitacional se genera un patrón de interferencia, finalmente se encuentra NANOGrav que utiliza los PTA's para detectar el paso por las GWs, este experimento fue el único que logró medir un fondo estocástico tal como está marcado con un círculo en la imagen.

    A su vez, en esta imagen podemos esquematizar la zona de interés para este trabajo que es la pintada de fucsia y representa a la dinámica de los campos de materia durante \emph{Reheating}. A su vez, se marca sobre ella las perspectivas de medición que tienen las EMRC como único instrumento optimista para la medición de este tipo de fondos.

    \section{EMRC}
      Como el único detector con expectativas de medición, entendamos rápidamente que son las EMRC. La idea de estas es que usan el efecto Gertsenshtein inverso para medir estas ondas, este efecto describe cómo es la interacción entre un campo electromagnético y las ondas gravitatorias. De este modo, se propone armar cavidades con un campo electromagnético estático, al pasar la onda genera un cambio en el volumen de las cavidades lo que genera perturbaciones en el campo interno y por la ley de Lenz se genera una corriente para contrarrestar este efecto, esta corriente tendrá una frecuencia semejante a la de la onda y si esta coincide con un modo de resonancia de la cavidad, la señal se amplifica y emite fotones. De este modo, la medición de estos fotones permite conocer las propiedades de las ondas. Y la medición de la autocorrelación o el uso de múltiples detectores permite utilizarlos no solo para medir ondas individuales sino también fonod cosmológicos.

  \section{Titulo Preheating Lineal}
    Empecemos con lo que trabajamos en esta tesis
  
  \section{Indice Preheating}
    En esta parte vamos a querer entender qué es la resonancia paramétrica y cómo esta afecta a la dinámica de los campos de materia cuando se tiene expansión del universo y cuando no.

  \section{Resonancia Paramétrica}
    Durante \textit{Preheating} mencionamos que el inflatón oscila alrededor del mínimo del potencial, este efecto transfiere la energía del inflatón a otros campos acoplados. A este fenómeno en que se transfiere energía se conoce como resonancia paramétrica y se puede describir a partir de la acción ... de la cual utilizando una métrica de fondo de tipo FLRW  permite extraer las siguientes ecuaciones dinámicas. Estas ecuaciones son diferenciales en derivadas parciales y están acopladas entre sí, por tanto, para poder resolverlas se hizo uso de un desarrollo perturbativo de estos dos campos y se estudiaron las perturbaciones del campo hijo acoplado.
  
  \section{Perturbaciones}
    Estas perturbaciones para el campo hijo se pueden escribir en el espacio de Fourier para poder trabajar con los modos de forma desacoplada, obteniendo una ecuación de la forma... siendo una ecuación del tipo oscilador armónico con disipación para cada modo y con una frecuencia de oscilación que depende del inflatón tiempo a tiempo

    Como primer estudio de estas ecuaciones vamos a considerar que no tenemos expansión del espacio-tiempo, lo que nos permite descartar el término de fricción con $H$ y escribir las ecuaciones en términos de un momento efectivo $\kappa$ y un parámetro que representa la intensidad efectiva del acople de los campos $q$.

    En particular, para el caso en que el potencial del inflatón es de forma cuadrática, se obtiene una ecuación conocida como de Mathieu mientras que para el caso de un potencial cuártico la ecuación se transforma en una ecuación del tipo Lamé. 
    
  \section{Análisis de Floquet}
    Tanto la ecuación de Mathieu como la de Lamé son ecuaciones del tipo oscilador/antioscilador y por tanto admiten soluciones de la forma ... De lo que se puede observar que existen valores de $\mu_k$ (conocido como coeficiente de Floquet) para los cuales se tiene una amplificación que son aquellos para los cuales $\Re{\mu} \neq 0$ si en cambio este coeficiente es imaginario puro sólo obtenemos por solución un seno o coseno que no nos daría amplificación. Esa misma idea se encuentra representada en esta figura donde se observan en negro la combinación de parámetros tal que no existe amplificación y en colores intensos los valores para los cuales la amplificación es mayor.
    
  \section{Amplificación por bandas}
    Si hacemos cortes transversales en la figura anterior se obtenienen cosas como estas dónde se observan las zonas en las que se tiene amplificación marcadas en naranja. El ancho de estas bandas al igual que la amplitud de las mismas nos dará el comportamiento del campo durante la resonancia, planteando dos regímenes conocidos como \textit{Narrow resonance} y \textit{Broad resonance}.
  
  \section{Tipos de resonancia: Narrow}
    El tipo de resonancia conocido como Narrow resonance o resonancia estrecha ocurre cuando el acoplamiento entre los campos es chico y en estos casos la amplitud del campo crece con el tiempo a la vez que oscila como se ve acá... y el mismo comportamiento se observa en el núimero de ocupación dónde se vé que este valor crece monótonamente con cada oscilación del inflatón.

  \section{Tipos de resonancia: Broad}
    El caso de resonancia ancha (o \textit{Broad resonance}) es el caso en que el acoplamiento de los campos es grande y por tanto se produce una mayor transferencia de energía pero al pasar el inflatón por el mínimo del inflatón se producen saltos abruptos en la amplitud del campo generando una dinámica escalonada para el número de ocupación del campo. En estos escalones cada salto se da como una ruptura de adiabaticidad generando un estallido de partículas.

  \section{Dinámica con expansión}
    Al nuevamente considerar la expansión del universo, uno puede definir variables comóviles para los campos de modo tal de extraer el término de fricción de la misma. Al hacer esto, la ecuación de perturbaciones para un potencial cuadrático se transforma en la siguiente... donde se puede observar que la frecuencia efectiva del oscilador armónico depende de la expansión del universo dando lo que se conoce como una resonancia estocástica.
    
    Para el caso del potencial cuártico la definición de estas variables comóviles recuperan las mismas ecuaciones que teníamos en el caso sin expansión permitiendo decir que la teoría es conforme.
    
    \section{Conexión a CosmoLattice}
      Hasta ahora estudiamos la dinámica de los campos en Recalentamiento a orden lineal y con eso surgen como preguntas: es posible estudiar los efectos no lineales que tienen estos campos? Es posible estudiar como esta dinámica puede llegar a producir GWs? La respuesta a esta pregunta es que sí, es posible, pero se requieren herramientas computacionales más complejas como las simulaciones de tipo Lattice como \CosmoLattice.

\end{document}